\section{SDE de preservación de varianza: VP-SDE}

\subsection{De DDPM a VP-SDE}

La distribución de transición hacia adelante de DDPM se define como:

$$
q(\mathbf{x}_t \mid \mathbf{x}_{t-1}) = \mathcal{N}\!\left(\mathbf{x}_t;\, \sqrt{1-\beta_t}\,\mathbf{x}_{t-1},\, \beta_t \mathbf{I}\right)
$$

La distribución de salto correspondiente es:

$$
q(\mathbf{x}_t \mid \mathbf{x}_0) = \mathcal{N}\!\left(\mathbf{x}_t;\, \sqrt{\bar{\alpha}_t}\,\mathbf{x}_0,\, (1 - \bar{\alpha}_t)\mathbf{I}\right)
$$

Donde $\bar{\alpha}_t = \prod_{s=1}^{t}(1-\beta_s)$.

\begin{warningbox}{Diferencia clave entre DDPM y SMLD}
SMLD solo aumenta la varianza (el media permanece inalterado), mientras que DDPM simultáneamente \textbf{reduce el media} y aumenta la varianza —el factor $\sqrt{1-\beta_t}$ hace que el media decaiga gradualmente hacia cero. Esto significa que en DDPM los puntos de datos están al mismo tiempo "arrastrados hacia el origen" y "perturbados por ruido", convergiendo finalmente a una distribución gaussiana estándar $\mathcal{N}(0, \mathbf{I})$.
\end{warningbox}

\subsection{Derivación de VP-SDE}

Definiendo la varianza continua $\beta(t)$ (asociada mediante $\beta_t^{\text{discrete}} = \beta(t) \cdot \Delta t$), la cadena de Markov discreta de DDPM se convierte en:

$$
\mathbf{x}_{t+\Delta t} = \sqrt{1 - \beta(t)\Delta t}\;\mathbf{x}_t + \sqrt{\beta(t)\Delta t}\;\boldsymbol{\epsilon}
$$

Utilizando la aproximación de Taylor de primer orden $\sqrt{1 - \beta(t)\Delta t} \approx 1 - \frac{1}{2}\beta(t)\Delta t$ y tomando el límite $\Delta t \to 0$:

$$
\mathrm{d}\mathbf{x} = -\frac{1}{2}\beta(t)\,\mathbf{x}\,\mathrm{d}t + \sqrt{\beta(t)}\,\mathrm{d}\mathbf{w}
$$

Esta es la \textbf{VP-SDE} (Stochastic Differential Equation de Preservación de Varianza). Los coeficientes correspondientes son:

\begin{itemize}
    \item \textbf{Coeficiente de arrastre}: $f(\mathbf{x}, t) = -\frac{1}{2}\beta(t)\,\mathbf{x}$ (arrastre lineal que atrae a $\mathbf{x}$ hacia el origen)
    \item \textbf{Coeficiente de difusión}: $g(t) = \sqrt{\beta(t)}$
\end{itemize}

\begin{importantbox}{Origen del nombre VP-SDE}
Si la covarianza de los datos iniciales es la matriz identidad $\mathbf{I}$ (es decir, $\mathrm{Cov}(\mathbf{x}_0) = \mathbf{I}$), entonces bajo VP-SDE, para todos los momentos intermedios $t$, la covarianza permanece igual a $\mathbf{I}$:

$$
\mathrm{Cov}(\mathbf{x}_t) = \mathbf{I}, \quad \forall\, t
$$

Esto es el significado de "preservación de varianza" (Variance Preserving). El término de arrastre $-\frac{1}{2}\beta(t)\mathbf{x}$ en la ecuación difusiva compensa exactamente la varianza inyectada por el término de difusión, manteniendo la varianza total constante.
\end{importantbox}

\subsection{$\alpha_t$ y $\sigma_t$ en VP-SDE}

Bajo el marco de Variational Diffusion Models, VP-SDE corresponde a:

$$
\alpha_t = e^{-\frac{1}{2}\int_0^t \beta(s)\,\mathrm{d}s}, \qquad \sigma_t^2 = 1 - \alpha_t^2 = 1 - e^{-\int_0^t \beta(s)\,\mathrm{d}s}
$$

\begin{knowledgebox}{Relación entre $\bar{\alpha}_t$ discreto y $\alpha_t$ continuo}
En DDPM discreto, $\bar{\alpha}_t = \prod_{s=1}^t (1-\beta_s)$, lo cual es una forma de producto. En el límite continuo, se convierte en una forma exponencial $\alpha_t = e^{-\frac{1}{2}\int_0^t \beta(s)\,\mathrm{d}s}$. Ambos son aproximaciones equivalentes: utilizando $\ln(1-x) \approx -x$ (cuando $x$ es muy pequeño), $\prod(1-\beta_s) \approx e^{-\sum \beta_s} \approx e^{-\int \beta(s)\,\mathrm{d}s}$. Sin embargo, ten en cuenta que el exponente en el producto es $-\sum \beta_s$, mientras que el exponente en la versión continua contiene un factor $\frac{1}{2}$; esta diferencia proviene de los detalles de las aproximaciones $\sqrt{1-\beta}$ y $1-\frac{\beta}{2}$.
\end{knowledgebox}

\subsection{Monotonía decreciente del SNR}

La relación señal-ruido (SNR) en VP-SDE es:

$$
\mathrm{SNR}(t) = \frac{\alpha_t^2}{\sigma_t^2} = \frac{e^{-\int_0^t \beta(s)\,\mathrm{d}s}}{1 - e^{-\int_0^t \beta(s)\,\mathrm{d}s}}
$$

Dado que $\beta(t) > 0$, la integral $\int_0^t \beta(s)\,\mathrm{d}s$ crece estrictamente con respecto a $t$; por lo tanto, el numerador decrece monótonicamente y el denominador crece monótonicamente, resultando en que el SNR decrece monótonicamente. Esto valida que DDPM es un caso particular de Variational Diffusion Models.

\begin{warningbox}{$\beta(t)$ debe ser positivo}
En DDPM y VP-SDE, $\beta(t) > 0$ es una restricción fundamental. Si $\beta(t)$ fuera negativo en algún momento, el SNR podría dejar de decrecer monótonicamente y el proceso hacia adelante ya no sería un proceso de "adición de ruido" válido. En implementaciones prácticas, $\beta_t$ se define típicamente como un programa lineal o coseno que garantiza ser siempre positivo.
\end{warningbox}

\subsection{Resumen de la sección}

VP-SDE es una generalización de DDPM al dominio del tiempo continuo: el término de arrastre atrae los datos hacia el origen, mientras que el término de difusión inyecta ruido; ambos están cuidadosamente equilibrados para mantener la varianza constante. A diferencia de VE-SDE, VP-SDE posee un arrastre no nulo.
